\documentclass[10pt,a4paper]{amsart}
\usepackage[margin=19mm]{geometry}
\usepackage{amsmath,amssymb,mathtools}
\usepackage[scr=rsfs,cal=euler]{mathalfa}
\usepackage{xfrac}
\usepackage[unicode,hidelinks,bookmarks=false]{hyperref}
\newcommand{\ie}{i.e.\ }
\newcommand{\cf}{cf.\ }
\newcommand{\resp}{respectively\ }
\newcommand{\Subsection}{\subsection}
\providecommand{\nobreakdash}{\nobreak\hbox{-}}
\setlength{\emergencystretch}{2em}
\multlinegap0pt
\usepackage{mathtools}
\usepackage{algorithm}\usepackage{algpseudocode}
\usepackage[shortlabels]{enumitem}
\usepackage{tcolorbox}
\newtheorem{lemma}{Lemma}
\newtheorem{theorem}{Theorem}
\newcommand{\EE}{\mathbb{E}}
\newcommand{\abs}[1]{\left\lvert #1 \right\rvert}
\newcommand{\cF}{\mathcal{F}}
\newcommand{\cN}{\mathcal{N}}
\newcommand{\cR}{\mathcal{R}}
\renewcommand{\frac}[2]{\tfrac{#1}{#2}}
\newcommand{\rbr}[1]{\left(#1\right)}
\newcommand{\sbr}[1]{\left[#1\right]}
\title{Sample Average Approximation for Distributionally Robust Optimization with $\phi$-divergences}
\author{Yan Li}
\date{Source: arXiv:2604.10855v4 (CC BY 4.0)}
\begin{document}
\maketitle

\noindent\emph{Source and licence.} Sample Average Approximation for Distributionally Robust Optimization with $\phi$-divergences. Version arXiv:2604.10855v4 (CC BY 4.0), distributed by arXiv under the Creative Commons Attribution 4.0 International licence, http://creativecommons.org/licenses/by/4.0/, as recorded in the arXiv licence metadata for this exact version and shown on the abstract page https://arxiv.org/abs/2604.10855v4. Full licence notice: \texttt{licenses/arxiv-2604.10855-CC-BY-4.0.txt}.\par\smallskip

\noindent\textit{Editorial scope.} Counted unit: the article's theorem thrm\_sample\_bernstein (source0.tex lines 1120--1135). The article's complete original proof of it is delivered with it (source0.tex lines 1138--1236, one \texttt{proof} environment of 99 lines). No mathematics is written, repaired or re-derived: the statement and the proof are the article's own lines, in the article's own notation, and no retained line is substituted or re-flowed.  Retained uncounted as the source-local context the proof needs: lines 804--818; lines 814--817; lines 849--855; lines 852--854; lines 925--930; lines 945--948; lines 1084--1096; lines 1088--1095.  Nothing was omitted from any retained span, so the package carries no omission record.  No retained line contains a citation, so the wrapper carries no bibliography and no bibliography entry.  No retained line contains a figure environment, an image-inclusion command or drawing code, so no image is delivered and the package declares no asset dependency.  Editorial formatting only, in addition: an AMS \texttt{amsart} preamble that restates the article's own declarations and macros used by the retained lines, namely the environments \texttt{lemma}, \texttt{theorem} on separate counters, together with the standard packages those lines need.  The retained environments print in this document as Lemma 1, Theorem 1 in order of appearance, whereas the article prints the same items with the numbering of its own full document.  The complete native source of the article as distributed in the arXiv e-print for this exact version is bundled unchanged as \texttt{sources/198296/source0.tex}.
\begin{lemma}\label{lemma_lipschitz}
For any divergence function $\phi$, any $L \geq 1$ and $\underline{\lambda} > 0$,
we have that
\begin{align}\label{def_f_truncate}
f_L(\lambda, \mu; \omega) \coloneqq  \lambda \tau + \mu + (\lambda \phi_L)^*(X(\omega) - \mu)
\end{align}
is Lipschitz continuous w.r.t $\mu$ (resp. $\lambda$) with modulus $L+1$ (resp. $\frac{2LB}{\underline{\lambda}} + \tau$) 
over domain 
$ [\underline{\lambda},  \frac{2B}{\tau}] \times [-B, B]$. 
In addition, we have 
\begin{align}\label{dual_stoch_obj_bound}
\abs{f_L(\lambda, \mu; \omega) } \leq (3 + 2L) B , ~ \forall 
\mu  \in [-B, B], ~ \lambda \in [\underline{\lambda},  \frac{2B}{\tau}].
\end{align}
\end{lemma}

\begin{align}
\abs{f_L(\lambda, \mu; \omega) } \leq (3 + 2L) B , ~ \forall 
\mu  \in [-B, B], ~ \lambda \in [\underline{\lambda},  \frac{2B}{\tau}].
\end{align}

\begin{lemma}\label{lemma_risk_via_truncation}
Suppose $\lim_{x \to \infty} \phi(x) / x = +\infty$.
For any $L \geq 1$, it holds that 
\begin{align}\label{ineq_truncate_tail_via_growth}
\cR(X) - \frac{B\tau}{g_\phi(L)} \leq \cR_L(X) \leq \cR(X).
\end{align}
\end{lemma}

\begin{align}
\cR(X) - \frac{B\tau}{g_\phi(L)} \leq \cR_L(X) \leq \cR(X).
\end{align}

 \begin{align}\label{eq_lip_m_constants}
 L(\epsilon)  =  g_\phi^{-1}\rbr{\frac{4 B\tau }{ \epsilon}} , ~~ \underline{\lambda}(\epsilon) =\frac{\epsilon}{4\tau}, ~ ~
 \overline{L} (\epsilon) = {L (\epsilon) +1 +  \frac{2L(\epsilon)B}{\underline{\lambda}(\epsilon)}  + \tau},  ~~ 
 M(\epsilon) = 
  (3 + 2L(\epsilon) ) B  .
 \end{align}

\begin{align}
\cR_L(X)   & \leq   \EE_P \sbr{f_L (\lambda^*, \mu^*; \omega)} \leq   \cR_L(X) + \underline{\lambda} \tau , \label{risk_via_emp_restrict} \\
   \cR_{n,L}(X)  & \leq  \EE_{P_n} \big[{f_L (\hat{\lambda}^*, \hat{\mu}^*; \omega)} \big] \leq \cR_{n,L}(X) +  \underline{\lambda}  \tau . \label{risk_via_emp_restrict_truncate}
\end{align}

\begin{lemma}\label{lemma_two_side_bernstein}
For any random variable $Z$ over $(\Omega, \cF, P)$ with $0 \leq Z \leq M$. 
Let $P_n$ denote the empirical distribution constructed from $n$ independent samples following distribution $P$. 
For any $\delta \in (0,1)$, with probability at least $1-\delta$, we have
\begin{align}
{ \EE_{P_n} \sbr{Z}  - \EE_P \sbr{Z} } & \leq  4  \bigg({
\sqrt{\frac{  { \EE_P [{Z^2}]} \log(3/\delta)}{n}} + \frac{M \log(3/\delta)}{n} 
} \bigg),  \label{eq_bernstein_1} \\
{  \EE_P \sbr{Z} - \EE_{P_n} \sbr{Z}   } & \leq  4 \bigg({
\sqrt{\frac{  { \EE_{P_n} [{Z^2}]} \log(3/\delta)}{n}} + \frac{M \log(3/\delta)}{n} 
} \bigg). \label{eq_bernstein_2}
\end{align} 
\end{lemma}

\begin{align}
{ \EE_{P_n} \sbr{Z}  - \EE_P \sbr{Z} } & \leq  4  \bigg({
\sqrt{\frac{  { \EE_P [{Z^2}]} \log(3/\delta)}{n}} + \frac{M \log(3/\delta)}{n} 
} \bigg),   \\
{  \EE_P \sbr{Z} - \EE_{P_n} \sbr{Z}   } & \leq  4 \bigg({
\sqrt{\frac{  { \EE_{P_n} [{Z^2}]} \log(3/\delta)}{n}} + \frac{M \log(3/\delta)}{n} 
} \bigg). 
\end{align} 

\begin{theorem}\label{thrm_sample_bernstein}
Suppose $\lim_{x \to \infty} \phi(x) / x = \infty$. 
For any $0 < \epsilon \leq B $, let 
$\underline{\lambda}(\epsilon), L(\epsilon)$, 
 $\overline{L}(\epsilon)$ and $M(\epsilon)$ be defined as in \eqref{eq_lip_m_constants}.
Then 
for any $\delta \in (0,1)$  
we have $\abs{\cR(X) - \cR_n(X)} \leq \epsilon$ with probability at least $1-\delta$
 if 
\begin{align*}
n \geq 
\frac{2048 B M(\epsilon)  \log (768 B^2 \overline{L}^2(\epsilon)/ (\epsilon^2 \tau \delta))}{\epsilon^2} 
+ 
\frac{32 M(\epsilon)  \log (768 B^2 \overline{L}^2(\epsilon)/ (\epsilon^2 \tau \delta))}{\epsilon} .
\end{align*}
\end{theorem}

\begin{proof}
We begin by noting that  
\begin{align}
%\label{pop_emp_min_loss_diff_2}
& \EE_{P_n} \big[{f_L (\hat{\lambda}^*, \hat{\mu}^*; \omega)}\big]
- 
\EE_{P} \sbr{f_L ({\lambda}^*, {\mu}^*; \omega)} \nonumber \\
\leq & 
\EE_{P_n} \sbr{f_L ({\lambda}^*, {\mu}^*; \omega)}
- 
\EE_{P} \sbr{f_L ({\lambda}^*, {\mu}^*; \omega)} 
\nonumber \\
\overset{(a)}{\leq} & 
 4 \bigg({
\sqrt{\frac{ M  {\EE_{P} \sbr{f_L ( \lambda^*,  \mu^*; \omega) }
} \log(12B^2/(\tau \varepsilon^2 \delta))}{n}} + \frac{ M \log(12B^2/(\tau \varepsilon^2 \delta))}{n} 
} \bigg)
\nonumber \\
\overset{(b)}{\leq} &   4 \rbr{
\sqrt{\frac{ M 
(B + \underline{\lambda} \tau  )
 \log(12B^2/(\tau \varepsilon^2 \delta))}{n}} + \frac{ M \log(12B^2/(\tau \varepsilon^2 \delta))}{n} 
}, \label{ineq_bernstein_at_pop_dual}
\end{align}
where $M = (3+2L)B$, $(a)$ follows from applying  \eqref{eq_bernstein_1} in  Lemma \ref{lemma_two_side_bernstein} while using  \eqref{dual_stoch_obj_bound} in Lemma \ref{lemma_lipschitz},
and $(b)$ follows from \eqref{risk_via_emp_restrict}.


For any $\varepsilon > 0$, let $\cN_1$ (resp. $\cN_2$) be an $\varepsilon$-net of $[\underline{\lambda}, \frac{2B}{\tau}]$ (resp. $[-B, B]$). 
In addition, for any $\lambda \in  [\underline{\lambda}, \frac{2B}{\tau}]$ and $\mu \in [-B, B]$, let 
$\lambda_\varepsilon$ and $\mu_\varepsilon$ be their corresponding closest points in $\cN_1$ and $\cN_2$ respectively. 
From \eqref{eq_bernstein_2} in Lemma \ref{lemma_two_side_bernstein} and \eqref{dual_stoch_obj_bound} in Lemma \ref{lemma_lipschitz}, we have that for any $\delta \in (0,1)$, with probability $1-\delta$,
\begin{align}\label{ineq_bernstein_eps_net}
{
\EE_{P} \sbr{f_L (\lambda_\varepsilon, \mu_\varepsilon; \omega)}  - \EE_{P_n} \sbr{f_L (\lambda_\varepsilon, \mu_\varepsilon; \omega)}
}
\leq 
 4 \rbr{
\sqrt{\frac{ M  {\EE_{P_n} \sbr{f_L (\lambda_\varepsilon, \mu_\varepsilon; \omega) }
} \log(12B^2/(\tau \varepsilon^2 \delta))}{n}} + \frac{ M \log(12B^2/(\tau \varepsilon^2 \delta))}{n} 
}
\end{align}
for any $ \lambda \in [\underline{\lambda}, \frac{2B}{\tau}], \mu \in [-B, B]$.
Consequently, we obtain
\begin{align}
&  \EE_{P} \sbr{f_L (\lambda^*, \mu^*; \omega)}
- 
\EE_{P_n} \big[{f_L (\hat{\lambda}^*, \hat{\mu}^*; \omega)} \big] \nonumber \\ 
 \leq &   \EE_{P} \big[{f_L (\hat \lambda^*, \hat \mu^*; \omega)} \big] 
-  \EE_{P_n} \big[{f_L (\hat{\lambda}^*, \hat{\mu}^*; \omega)} \big]  \nonumber  \\
 \overset{(c)}{\leq} &   \EE_{P} \big[{f_L (\hat{\lambda}^*_\varepsilon, \hat{\mu}^*_\varepsilon; \omega)}\big] 
-  \EE_{P_n} \big[{f_L (\hat{\lambda}^*_\varepsilon, \hat{\mu}^*_\varepsilon; \omega)} \big] + 2 \overline{L} \varepsilon \nonumber \\
 \overset{(d)}{\leq}  &  
 4 \bigg({
\sqrt{\frac{ M  {\EE_{P_n} \sbr{f_L (\hat \lambda^*_\varepsilon, \hat \mu^*_\varepsilon; \omega) }
} \log(12B^2/(\tau \varepsilon^2 \delta))}{n}} + \frac{ M \log(12B^2/(\tau \varepsilon^2 \delta))}{n} 
} \bigg) + 2 \overline{L} \varepsilon \nonumber  \\
 \overset{(e)}{\leq}  &   
  4 \rbr{
\sqrt{\frac{ M 
(B + \underline{\lambda} \tau + \overline{L} \varepsilon)
 \log(12B^2/(\tau \varepsilon^2 \delta))}{n}} + \frac{ M \log(12B^2/(\tau \varepsilon^2 \delta))}{n} 
}  + 2 \overline{L} \varepsilon, 
\label{pop_emp_min_loss_diff_1}
\end{align}
where 
$(c)$ follows from Lemma \ref{lemma_lipschitz} and the definition of $\overline{L} =  {L+1 +  \frac{2LB}{\underline{\lambda}} + \tau}$, $(d)$ follows from \eqref{ineq_bernstein_eps_net}, and $(e)$ follows from \eqref{risk_via_emp_restrict_truncate} and Lemma \ref{lemma_lipschitz}.

Now combining \eqref{risk_via_emp_restrict}, \eqref{risk_via_emp_restrict_truncate}, \eqref{ineq_bernstein_at_pop_dual}, and \eqref{pop_emp_min_loss_diff_1}, we obtain 
\begin{align*}
\cR_{n,L}(X) & \leq \cR_L(X) + \underline{\lambda} \tau + 4 \rbr{
\sqrt{\frac{ M 
(B + \underline{\lambda} \tau  )
 \log(12B^2/(\tau \varepsilon^2 \delta))}{n}} + \frac{ M \log(12B^2/(\tau \varepsilon^2 \delta))}{n} 
}, \\
\cR_L(X) & \leq \cR_{n,L}(X) + \underline{\lambda} \tau 
+  4 \rbr{
\sqrt{\frac{ M 
(B + \underline{\lambda} \tau + \overline{L} \varepsilon)
 \log(12B^2/(\tau \varepsilon^2 \delta))}{n}} + \frac{ M \log(12B^2/(\tau \varepsilon^2 \delta))}{n} 
}  + 2 \overline{L} \varepsilon,
\end{align*}
which together with \eqref{ineq_truncate_tail_via_growth} in Lemma \ref{lemma_risk_via_truncation}, implies
\begin{align*}
\abs{ \cR_{n} (X) - \cR(X)} 
\leq 
 \underline{\lambda} \tau 
+  4 \rbr{
\sqrt{\frac{ M 
(B + \underline{\lambda} \tau + \overline{L} \varepsilon)
 \log(12B^2/(\tau \varepsilon^2 \delta))}{n}} + \frac{ M \log(12B^2/(\tau \varepsilon^2 \delta))}{n} 
}  + 2 \overline{L} \varepsilon
+ \frac{B\tau}{g_\phi(L)}.
\end{align*}
The desired claim then follows by substituting the choice of 
$\underline{\lambda} = \underline{\lambda}(\epsilon)$,
$L = L(\epsilon)$, and $\varepsilon = \frac{\epsilon}{8 \overline{L}(\epsilon)}$ into the above relation, 
while noting that $\epsilon \leq B$.
\end{proof}

\end{document}
